\section{Local rules and the scalar theorem}\label{sec:intro}
The fourth paper proves a scalar linear storage bound on triangle-free
subcubic graphs, including a more general directed-support form and an
exact matching-number consequence~\cite{BrendeCurved}. A scalar row can
solve for any coordinate with a nonzero coefficient. We isolate this
property and retain the support combinatorics, without assuming scalar
linearity. The resulting generating set applies to every alphabet and
compatible code on the same supports.

Closure operators and their entropy are an established framework for
network coding~\cite{Gadouleau2014}. Our use of closure and polymatroids
is within that framework. The result here is the matching-size generator
bound for the specified local rules, together with a schedule that charges
each center once. We do not claim a bound for arbitrary nonlinear recovery
on all triangle-free subcubic graphs.

Sections~\ref{sec:rank-one} and~\ref{sec:higher-rank} treat a different
linear class: arbitrary coordinate widths with at most one higher-rank
component per row. These bounds measure actual coordinate images and
charge exceptional rows by their exact quotient deficit.
Sections~\ref{sec:squares} and~\ref{sec:quadratic} treat arbitrary block
ranks through the actual two-step defect. They give a weighted graph
bound, a paid kernel bound, and an exact quadratic rank identity.

Every numbered result below has a Lean proof. Those declarations are the
proofs; the English outlines are explanatory guides. The formalization
imports the fourth paper's library and reuses its good-chain geometry and
Tutte--Berge machinery. It proves new generating-set statements rather
than assuming the old scalar dimension inequality for nonlinear objects.

Let $V$ be finite. At each $v$ choose $N_v\subseteq V\setminus\{v\}$ and
write $R_v=\{v\}\cup N_v$. The general directed hypotheses are
\begin{align}
 &p,q\in N_v\ \Longrightarrow\ q\notin N_p,
   \label{eq:no-triangle}\\
 &|N_v|\le2\quad\text{or}\quad
   \bigl(|N_v|=3\ \text{and}\ v\in N_u\Longrightarrow u\in N_v
       \text{ for every }u\bigr).
   \label{eq:star}
\end{align}
They are exactly the no-recovery-triangle and reciprocity conditions used
in the original directed theorem. In particular, they hold whenever
$N_v\subseteq N_G(v)$ for a finite simple triangle-free graph $G$ of maximum
degree at most three. In the directed statements, $G$ may be any simple
undirected graph containing every pair $vu$ with $u\in N_v$.
Write $\nu(G)$ for its maximum matching size.

In Lean a matrix records the supports $N_v$. Only its zero pattern is
used in the generator proof. Arbitrary named sets are encoded by a zero-one
rational matrix; this introduces no field assumption on the recovered data.

\section{A fixed generating set}\label{sec:generators}
Starting with known coordinates $B\subseteq W$, a rule centered at
$v\in W$ may mark $x\in R_v\cap W$ known whenever all other members of
$R_v\cap W$ are already known. A set generates $W$ if finitely many such
steps mark every member of $W$ known. These are combinatorial rules;
restricting a row to $W$ does not assert that an arbitrary projected code
can recover without its omitted coordinates.

\begin{theorem}[Half-size generators and matching-size generators]
\label{thm:generators}
\formal{ReversibleRules.lean}{exists_generators_half, exists_generators_matching}
Under \cref{eq:no-triangle,eq:star}, every $W\subseteq V$ has a generating
set $B\subseteq W$ with $|B|\le\lfloor|W|/2\rfloor$. On $W=V$, if $G$
contains the prescribed pairs, there is such a $B$ with $|B|\le\nu(G)$.
The choices depend only on the row supports.
\end{theorem}
\begin{outline}
Induct on $|W|$. If $|W|$ is positive and even, choose one seed and apply
the odd-size induction bound to the remaining set, supplying the removed
seed whenever a truncated row needs it. For odd $|W|$, the original
small-recovery lemma supplies a root with at most two parents in $W$.
A good chain starts at this root. When a known center has two parents
outside the chain, add both and seed one; its row derives the other.
After $j$ additions the chain has $2j+1$ vertices and at most $j+1$ seeds.
Reciprocity ensures that every chain vertex has at most two parents outside.

If a chain row has just one outside parent, derive that parent without a
seed. The enlarged set has at most half as many seeds as vertices; apply
induction to its complement. If a pair remains, extend the chain. Finally,
if no row has outside parents, either the singleton root derives itself
without a seed, or the last pair shares a parent and has no prescribed
edge between its members by \cref{eq:no-triangle}. One member of that last
pair is then derivable from the preceding chain, so its seed can be removed.
Again the chain and complement meet the half-size bound. Finiteness
terminates the chain construction.

For the matching bound, use the original library's Tutte--Berge partition:
choose $U\subseteq V$ with
$\nu(G)=|U|+\sum_A\lfloor|A|/2\rfloor$, where $A$ ranges over the
components of $G-U$. Seed $U$ and a half-size generating set in each
component. Every omitted premise of a component row is in $U$, since no
graph edge joins different components. Thus the component derivations lift
to the whole vertex set, with at most $\nu(G)$ seeds.
\end{outline}

\begin{theorem}[Productive schedule]\label{thm:schedule}
\formal{Schedules.lean}{Generates.exists_schedule, Schedule.centers_nodup, exists_matching_schedule}
The set $B$ in the matching conclusion can be accompanied by a finite list
of steps $(v_i,x_i)$, chosen from the supports, which marks only previously
unknown targets and generates $V$. No center $v_i$ occurs twice.
\end{theorem}
\begin{outline}
If a derivable target remains unknown, descend through its finite derivation
until finding a rule all of whose other premises are already known. Apply
that productive rule. The unknown set strictly shrinks. After using a row,
all of its members are known, so that center can never again produce a new
target. This last argument requires no graph hypothesis.
\end{outline}

A code $\mathcal C\subseteq\mathcal A^V$ is \emph{reversible on these rows}
if, for every $v$ and every $x\in R_v$, agreement of two codewords on
$R_v\setminus\{x\}$ implies agreement at $x$. Ordinary recovery of $v$
from $N_v$ is only one of these directions.

\begin{theorem}[Alphabet-independent information set]\label{thm:information}
\formal{ReversibleRules.lean}{reversible_information_set, named_rules_information_set, reversible_card_le, reversible_finrank_le}
There is a set $B$ of size at most $\nu(G)$, determined by the supports,
such that projection onto $B$ is injective on every code reversible on the
rows. In particular, for a finite nonempty alphabet of size $q$,
\[
 |\mathcal C|\le q^{\nu(G)}.
\]
If the coordinates lie in one finite-dimensional vector space $A$ over any
field $K$, and $\mathcal C\le A^V$ is linear and reversible, then
\[
 \dim_K\mathcal C\le\nu(G)\dim_K A.
\]
\end{theorem}
\begin{outline}
Two codewords agreeing on $B$ agree at every step of a generating derivation.
Projection onto $B$ is therefore injective. Count the possible records in
the finite-alphabet case. In the linear case this projection is an injective
linear map into $A^B$, whose dimension is $|B|\dim_K A$.
\end{outline}
The set $B$ is an information set, not an efficient native tape decoder.
The theorem gives no bound on the cost of evaluating an individual inverse.

\section{Linear-block consequences and the original theorem}\label{sec:blocks}

\begin{theorem}[Injective participating blocks]\label{thm:blocks}
\formal{BlockRecovery.lean}{injective_block_information_set, injective_block_finrank_le}
Let $G$ be finite, simple, triangle-free, and subcubic. Let $A$ be a
finite-dimensional vector space over any field $K$, and let
$T_{vu}:A\to A$ be linear maps with $T_{vv}=I$ and $T_{vu}=0$ off the
edges for $v\ne u$. Suppose every nonzero $T_{vu}$ is injective.
For every linear subspace $\mathcal C\le A^V$ satisfying
$\sum_uT_{vu}x_u=0$ at every $v$,
\[
 \dim_K\mathcal C\le\nu(G)\dim_K A.
\]
When $A$ has positive dimension, the corresponding information set is chosen
from the nonzero block pattern and works for every, possibly nonlinear,
solution set. In finite dimension the nonzero injective endomorphisms are
invertible; no commutation assumption is imposed.
\end{theorem}
\begin{outline}
Subtract two solutions agreeing at all but one participating member of a
row. The remaining block applied to their target difference is zero;
injectivity kills that difference. Thus every participating row is
reversible, and \cref{thm:information} applies. The zero-dimensional case
has zero code dimension directly.
\end{outline}

\begin{corollary}[Original scalar statements]\label{cor:scalar}
\formal{ReversibleRules.lean}{original_subCode_half, original_matching_bound}
Let $M$ be a scalar matrix with unit diagonal, prescribed directed supports
$N_v$, and \cref{eq:no-triangle,eq:star}. Its original restricted code
\[
 \mathcal C_W=\{x:x_u=0\ (u\notin W),\quad (Mx)_v=0\ (v\in W)\}
\]
satisfies $\dim\mathcal C_W\le\lfloor|W|/2\rfloor$ whenever
\cref{eq:star} holds at the centers in $W$. If a graph $G$ contains every
prescribed pair, then $\dim\ker M\le\nu(G)$.
\end{corollary}
\begin{outline}
A scalar row solves for each variable with nonzero coefficient. This
remains true after outside coordinates are set to zero. Use the generating
set of \cref{thm:generators} and inject its restricted code into the scalar
seed coordinates. The full-code matching statement follows from
\cref{thm:information}. Thus both statements are derived from the stronger
rules theorem, not from a second copy of the original scalar induction.
\end{outline}

\begin{corollary}[Exact scalar minrank]\label{cor:minrank}
\formal{BlockRecovery.lean}{original_graph_matching, original_minrank}
For a triangle-free subcubic graph $G$, the minimum rank, over any field,
of a scalar unit-diagonal matrix with off-diagonal support on $G$ is
$|V|-\nu(G)$.
\end{corollary}
\begin{outline}
The preceding nullity bound and rank--nullity give the lower bound. Reuse
the original library's matching-supported construction for attainment,
and apply the new upper nullity bound to get equality. This is a
rederivation of the original minrank result, not a new minrank value.
\end{outline}
